% =============================================================================
% Cover Letter for the NL2Semantic2SQL v2 submission.
% Revised for v2: reports BIRD significance (p=0.0106), cross-lingual 100q,
% Robustness 40q expansion, 6-way ablation, and honest OOM gap disclosure.
% =============================================================================
\documentclass[11pt]{letter}
\usepackage[margin=1in]{geometry}
\usepackage{times}
\usepackage{url}
\usepackage[hidelinks]{hyperref}

\signature{Zhou Ning \\ \textit{Beijing SuperMap Software Co., Ltd.}}
\address{Zhou Ning \\ Beijing SuperMap Software Co., Ltd. \\ \texttt{zhouning1@supermap.com}}

\begin{document}

\begin{letter}{The Editor-in-Chief \\ \emph{International Journal of Geographical Information Science}}

\opening{Dear Editor,}

We are pleased to submit a substantially revised version of our manuscript, ``\textbf{NL2Semantic2SQL: A Cross-Domain Framework for Natural Language to SQL over Geospatial and Enterprise Data Warehouses},'' for consideration as a regular research article in \emph{IJGIS}. This submission addresses the reviewer feedback on the previous version and contributes new experimental evidence not present in the earlier draft.

\medskip
\noindent\textbf{What is new in this version.}

\begin{enumerate}
\item \textbf{BIRD statistical significance achieved.} The previous version reported a directional but non-significant BIRD result (EX $+0.027$, McNemar $p=0.136$) on 495 questions. Through error-attribution analysis of round-1 failures, we identified three recurring gap patterns (missing \texttt{DISTINCT} on one-to-many joins, over-JOIN errors when a single table suffices, string-concatenation vs.\ separate-column output format) and added three targeted rule sections to the grounding prompt. On 108 paired questions, the round-2 pipeline achieves EX $=0.593$ vs.\ baseline $0.500$ ($+0.093$, McNemar $p=0.0106$, \textbf{statistically significant at $\alpha=0.05$}). The Round-1 $\to$ Round-2 transition ($p=0.2266 \to p=0.0106$) on comparable sample sizes isolates the causal contribution of the new rule sections, rather than attributing the gain to sample-size inflation.

\item \textbf{Cross-lingual evaluation on 100 paired questions.} The previous version contained only a preliminary stress test with four-percentage-point loss. We now report 100 question pairs: 50 Chinese-translated BIRD questions (EX $=0.560$ vs.\ English same-question EX $=0.660$, bounded degradation of $-0.100$ attributable to value translation drift, not SQL generation failure --- validity remains $0.980$); and 50 native-Chinese GIS questions (EX $=0.820$). This is, to our knowledge, the first execution-based cross-lingual text-to-SQL study that combines geospatial and warehouse tracks.

\item \textbf{Robustness expanded to 40 questions across 6 categories}, including a new Schema Hallucination category that probes behavior on non-existent tables and columns. The pipeline achieves $1.000$ Success Rate on five categories (Security Rejection, Anti-Illusion, Schema Hallucination, Schema Enforcement, Data Tampering Prevention), and we \textbf{honestly disclose} a $0/8$ OOM Prevention gap: the intent-based \texttt{LIMIT} injection does not fire on naturally phrased ``show me all $X$'' queries that classify as \texttt{ATTRIBUTE\_FILTER} rather than \texttt{PREVIEW\_LISTING}. We name this as future work (an orthogonal large-table guard layer) rather than hide it.

\item \textbf{Six-way ablation with domain-scoping finding.} We conduct a full 6-way ablation on the GIS 125-question benchmark. An important negative result emerges: the round-2 rule sections and warehouse join hints do \emph{not} aid spatial accuracy (all deltas within $\pm 0.016$), confirming that these rules are domain-scoped to warehouse queries. This motivates future work on domain-selective grounding --- enable warehouse rules only when the intent router detects a warehouse query --- which we note as a concrete next step.

\item \textbf{Expanded GIS benchmark to 125 questions} (85 Spatial + 40 Robustness), from the original 20-question pilot and 100-question intermediate version.
\end{enumerate}

\medskip
\noindent\textbf{Contributions restated against the current state of the literature.}
\begin{itemize}
\item Unified semantic-to-SQL pipeline for both GIS (PostGIS) and enterprise warehouses (BIRD), with error-attribution-driven grounding rule augmentation that produces statistically significant gains on \emph{both} tracks --- addressing a gap in Spider, BIRD, Spider~2.0 and BEAVER, none of which cover executable spatial SQL.
\item First execution-based cross-lingual text-to-SQL evaluation spanning both GIS and warehouse domains, with bounded degradation analysis.
\item Honest reporting of domain-scoped rule contributions (GIS 6-way ablation) and the OOM Prevention limitation, with concrete future-work paths.
\end{itemize}

\medskip
\noindent\textbf{Reproducibility.} All 125 GIS benchmark questions with gold SQL, the 108-question BIRD R2 evaluation records, the 100-question cross-lingual pairs, the 6-way ablation runner, the DIN-SQL comparison scripts, and per-question execution logs are released in the project repository. Each table in the manuscript corresponds to a specific run directory under \texttt{nl2sql\_eval\_results/}. The releasing commits are identified in the Reproducibility section.

\medskip
\noindent\textbf{Originality and conflicts of interest.} The manuscript is original work and has not been published or submitted elsewhere. No portion has appeared in any conference or workshop. The authors declare no competing interests.

\medskip
\noindent\textbf{Revision disclosure.} A previous version of this manuscript underwent external review; this revision responds to that feedback via the accompanying ``Response to Reviewers'' document, which details how each comment is addressed in the new version.

We thank you for considering our submission and look forward to the reviewers' feedback.

\closing{Sincerely,}

\end{letter}

\end{document}
